\subsection{RL-based Legged Locomotion}
Reinforcement learning has seen significant progress in quadrupedal locomotion control, largely driven by advances in physics simulators like Mujoco~\cite{todorov2012mujoco} and IsaacGym~\cite{makoviychuk2021isaac}. This high-fidelity, GPU-accelerated platform allows for efficient RL training by simulating complex environments~\cite{rudin2022learningwalkminutesusing}.  

To mitigate the uncertainty caused by partial observation, the teacher-student framework~\cite{kumar2021rma, doi:10.1126/scirobotics.abc5986} was proposed, enabling quadrupedal robots to adapt to unknown terrains and greatly improving locomotion flexibility and robustness. Much of the work~\cite{pmlr-v205-agarwal23a, Cheng2024parkour, kumar2022adapting} has continued in this structure with impressive results.
Walk-These-Way~\cite{margolis2023walk} proposes a method to improve the generalization of quadrupedal locomotion by training the robot in diverse environments.
Another attempt is introducing environmental representation, such as~\cite{ji2022concurrent, nahrendra2023dreamwaq, long2024him}.
DreamWaQ~\cite{nahrendra2023dreamwaq}, leverages variational autoencoders (VAE)~\cite{higgins2016beta} to learn environmental representations, while HIMloco~\cite{long2024him} employs contrastive learning techniques to learn terrain representations.
Recent advances~\cite{lyu2023composite, lyu2024rl2ac} introduce model-based control~\cite{kim2019highly} and adaptive control~\cite{ioannou1996robust} to improve the interpretability and robustness of RL-based locomotion.
To address the challenges of the sim2real gap, domain randomization has been commonly used~\cite{xie2021dynamics, tan2018sim, kumar2021rma, nahrendra2023dreamwaq, long2024him}.
It is a class of methods in which the policy is trained with a wide range of environment parameters and sensor noise to learn behaviors that are robust in this range. However, domain randomization trades optimality for robustness, sometimes leading to an over-conservative policy.

\subsection{Model-based RL}
Model-based Reinforcement Learning (MBRL) incorporates a learned dynamics model of the environment, which predicts future states given the current state and action. This allows MBRL methods to plan and use collected data efficiently. Popular frameworks include Dreamer series~\cite{hafner2019dream, hafner2020mastering, hafner2023mastering}, which uses a learned world model to perform planning in latent space, and TDMPC series~\cite{hansen2022temporal, hansen2023td}, which integrates model predictive control to improve sample efficiency.
In the legged locomotion tasks, recent MBRL approaches have demonstrated the power of learned dynamics models. DayDreamer~\cite{wu2023daydreamer}, inspired by Dreamer, enables effective planning and control for robotic locomotion. PIP-loco~\cite{shirwatkar2024pip} also adopts a model-based approach and combines it with model predictive control, resembling TDMPC. These methods utilize the learned model for trajectory planning or representation learning, leveraging the learned dynamics model to enhance performance.
In contrast to these works, our approach utilizes the dynamics model not for planning or representation, but as a reference to predict the next state, thereby more directly affecting policy learning.